% !PRINTED 393
\weberpage{393}
This case therefore always leads to a real double tangent.

But if, \gesperrt{secondly}, the pair of complexes 1), 2) is azygetic, then
every pair of the one complex contains an element that also occurs in the
other complex, and an element that does not occur in the other. If, then,
say $\eta$ occurs in the complex 2), we may, without impairing generality,
assume $\eta = \eta_1'$, $\eta' = \eta_1$,
and we obtain the following pairs in 1) and 2)
\weq{2}{
\begin{array}{rl}
1) & \xi\eta,\ \xi_1\eta' \\
2) & \xi_1'\eta,\ \xi'\eta'.
\end{array}
}
If, further, $\eta_2$ occurs in 1) and 2), then $\eta_2$ cannot be equal to
$\xi_2'$, since otherwise 1) would contain the pair $\xi_2\xi_2'$ and would
consequently be real.

If $\eta_2 = \eta_2'$, then $\eta_2$ is \gesperrt{real}. But if $\eta_2$ is equal
to an element of the three last pairs of 2), then we may without
restriction assume it $= \eta_3'$, hence $\eta_2 = \eta_3'$, $\eta_3 = \eta_2'$.

But then, by the theorem of \S.\,96, 3., both $\eta, \eta_2, \eta_2'$ and
$\eta, \eta', \eta_2$ are azygetic. The two pairs $\eta\eta_2$, $\eta'\eta_2'$
determine two complexes, of which the second contains neither $\eta$ nor $\eta_2$,
and which therefore form a \gesperrt{syzygetic pair}. With this we are led back
to the first assumption, and our theorem 1. is proved.

There are therefore always at least \gesperrt{two real} double tangents.

\gesperrt{2. There is always at least one system of four real
syzygetic double tangents.}

In the proof of this theorem we start from a real pair $x\,y$, which
by 1. always exists, and consider the real complex that is determined
by this pair. Since the conjugate double pairs are two pairs each,
there must, among the six pairs of this complex,
either occur a second real pair, and then theorem
2. holds, or it must contain a conjugate pair $\xi\xi'$.

In the latter case we consider the syzygetic triple of complexes
determined by these two pairs, in which all 28 double tangents
occur (\S.\,97, 4.):
\[
\begin{array}{rl}
1) & x\,y,\ \xi\xi' \\
2) & x\xi,\ y\xi',\ \xi_1\eta_1 \\
3) & x\xi',\ y\xi,\ \xi_1'\eta_1',
\end{array}
\]
%CONT>>
